\Needspace{9\baselineskip}\pdfbookmark[1]{Göttingische gelehrte Anzeigen. 1811 May 25.}{prasse-1811}
\begin{center}
\rule{0.60\linewidth}{0.3pt}\\
Göttingische gelehrte Anzeigen. 1811 May 25.\\
\rule{0.60\linewidth}{0.3pt}
\end{center}

\textit{Logarithmische Tafeln für die Zahlen, Sinus und Tangenten, neu geordnet
von \textup{\textsc{Moritz von Prasse}}, ordentlichem Professor der Mathematik zu Leipzig.} 80 pages in octavo. Leipzig. Distributed by P. J. Besson.

These tables contain the same material as \textsc{Lalande}'s popular small tables: the logarithms of all numbers up to 10000, and the logarith\SourceLeaf{251}%
ms of the sines and tangents for every single minute of the quadrant, all to five decimal places. But Mr. von \textsc{Prasse} has reduced this to one third of the number of pages, in a format not much larger, by adopting the arrangement usual in larger tables: ten logarithms in each line, with the initial digits given only once. All the differences have had to be omitted, however. Thus what is gained in compactness seems to be lost again in convenience. Since this can be judged only from actual use, the reviewer, who was accustomed to \textsc{Lalande}'s small tables, set them aside for a time and tried using the present ones. He found that, for the logarithms of numbers, which occupy 31 pages here and 111 in \textsc{Lalande}, the advantage of having to turn far fewer pages appreciably outweighs those minor inconveniences, and he therefore likes using these new tables. He did not find this so with the trigonometric tables, which occupy 40 pages here and 90 in \textsc{Lalande}, particularly because Mr. von \textsc{Prasse} lets the arcs run from 0 to 90 degrees, thus separating the sines and tangents from the cosines and cotangents. This is all the more troublesome because one very often needs, for example, the sine and cosine of an arc together, or wants the cosine from the sine or tangent without needing the arc itself. Here he would therefore always prefer \textsc{Lalande}'s tables, and he wishes that Mr. von \textsc{Prasse} had instead divided each page once more across the middle with a horizontal line. This would have avoided that unwelcome separation while still allowing the material to be compressed into the small space.

These tables also differ in that the last digit of each logarithm, whenever it has been increased, is set in a different typeface. Mr. von P. believes that this can promote greater accuracy in calculations. But since in practice one mostly has to calculate with logarithms that must be interpolated, taking this circumstance into account is troublesome. Whenever one believes that logarithms reliable only to half a unit in the fifth decimal place cannot give sufficiently accurate results, it is better to use larger tables with six or seven decimal places. The reviewer therefore cannot regard this arrangement, which, generally speaking, can indeed serve to \SourceLeaf{252}
\emph{double} the accuracy of the calculation, as very useful, especially since the italic figures beside the others are unpleasant to the eye and here and there are not sharp enough. Otherwise the printing is neat; only those who use such tables a great deal will wish for heavier paper.

